\section{Purpose of this chapter}
\label{sec:intro}

This chapter reports the Phase~2 results of the climate-firestorm simulation.
The purpose is to describe, under documented limits, how eight network topologies, homogeneous versus heterogeneous identity mixes, and two auditor instruments move misinformation scores on six geoengineering articles.
The chapter is a thesis-style results account.
It is not a reprint of the CIKM linear-chain study on commercial and crime news \parencite{maurya2025simulating}, and it is not a Twitter digital twin of Debnath and colleagues' 814,924-tweet observational map \parencite{debnath2023conspiracy}.

The theoretical spine remains online firestorms as defined by \textcite{pfeffer2014firestorms}: a sudden discharge of large quantities of negative word-of-mouth in social media, often with high affective charge.
Their Outlook section lists \emph{seven} interrelated factors, not six.
Section~\ref{sec:dnet} restates those seven factors from the Debnath-versus-D-net mapping and records the thesis six-knob list as a computational remapping.
Cross-media dynamics has no engine knob and is held.

Phase~2 asked a narrower empirical question than a full Pfeffer confirmation.
Holding surprise as a drip seed and holding the clock at eight ticks, do topology, identity mix, and auditor mode change live-cell mean \MI{} on climate seeds?
The harvest answers that question only descriptively.
\texttt{summary.json} sets \texttt{thesisGrade} to false.
$N=1$ with graph seed 42 forbids error bars and replicate inference.

The remainder of the chapter follows a single path.
Section~\ref{sec:methods} snapshots the protocol.
Section~\ref{sec:instruments} separates discrete from continuous \MPR{}.
Section~\ref{sec:topology} reports all eight topologies.
Section~\ref{sec:hhe} reports personas and mixes without averaging the two instruments.
Section~\ref{sec:designs} reports the five numbered comparisons from the full analysis (C1--C5).
Section~\ref{sec:pooled} states that cell pooling concatenates rows and is not a physical super-graph.
Section~\ref{sec:dnet} compares D-net structure and Pfeffer observables with the Debnath hashtag fallback.
Section~\ref{sec:limits} states limitations.
Section~\ref{sec:conclusion} closes.

Figures are copies of \path{thesisExperiment/analysis_full/figures/} (\texttt{fig00}--\texttt{fig22}).
Comparison panels follow \texttt{COMPARISONS.md}.
Exploratory rank tests in that file are not treated as replicate inference.
